Processos Pontuais Temporais (\gls{TPP}) modelam fluxos de eventos em tempo contínuo, e os TPPs baseados em Transformers se tornaram uma ferramenta forte para essa tarefa. Na prática, muitas vezes precisamos treinar em sequências curtas e depois rodar o modelo em sequências bem mais longas, e é exatamente aí que esses modelos falham: sua acurácia colapsa quando as sequências de teste são mais longas do que as vistas no treinamento. Neste trabalho, propomos o HoTHP (Hyperbolic Rotary Position Embedding-based Transformer Hawkes Process), um modelo Transformer para processos de Hawkes que incorpora um viés de recência diretamente no kernel de atenção. Em vez do kernel trigonométrico oscilante do RoTHP (Rotary Position Embedding-based Transformer Hawkes Process), o HoTHP usa um kernel hiperbólico que decai de forma suave e monotônica com o tempo entre os eventos, de modo que o modelo se mantém previsível mesmo em distâncias que nunca viu no treinamento. O decaimento é definido por um único parâmetro aprendido e é monotônico por construção. Como esse viés atua na atenção e não na função de intensidade, o HoTHP ainda modela tanto processos auto-excitatórios quanto auto-inibitórios. Avaliamos o modelo sob o protocolo train-short, test-long, com sequências de teste até $5\times$ mais longas que o treinamento, em dados sintéticos de Hawkes e em quatro datasets reais (Amazon, Stackoverflow, Taxi e Retweet). Nos dados sintéticos, a negativa do log da verossimilhança (\gls{NLL}) do HoTHP se mantém estável à medida que as sequências crescem, enquanto o RoTHP degrada. Nos dados reais, o HoTHP obtém a melhor NLL em extrapolação entre os modelos Transformer, e no Retweet a diferença é drástica ($-0{,}651$ contra $47{,}6$ do RoTHP, após normalização temporal). Seu ganho mais claro está na acurácia de predição de tipo: o HoTHP é o modelo mais acurado no Amazon e no Stackoverflow, à frente até do forte baseline recorrente Neural Hawkes Process (NHP). Seu custo adicional de treinamento ($1{,}05\times$ a $2{,}24\times$) é pago uma única vez e garante inferência estável em sequências bem mais longas.

% Separe as palavras-chave por ponto
\palavraschave{Processos Pontuais Temporais. Processo de Hawkes. Transformers. Extrapolação de comprimento.}
